As a final application, we optimize the learning rate schedule
of stochastic gradient descent (SGD) for training a CIFAR-10 classifier.
By following the metagradients with respect to the learning rate
at each step of training, our procedure matches grid searching
over standard learning rate schedules---despite starting with na\"ive
hyperparameters (a flat learning rate).

Unlike the other applications discussed here,
metagradients do not unlock state-of-the-art performance.
Instead, we discuss this application to illustrate the flexibility
of \replay{}, and in particular its ability to optimize metaparameters
that do not directly affect the loss landscape (i.e., that only affect
the model via the optimization trajectory).
As we discuss in Section \ref{sec:related},
approximate metagradient estimators cannot apply to these metaparameters.

\subsubsection{Setting}
To put learning rate schedule optimization into the metagradient framework,
we parameterize a schedule as a vector $\bm{\eta} \in \mathbb{R}^k$
comprising $k$ evenly-spaced keypoints,
so that the learning rate at iteration $t$ is given by
\begin{equation}
  \label{eq:lr_schedule}
  \eta(t) = \eta_{\lfloor kt/T \rfloor} + \frac{kt/T - \lfloor kt/T
  \rfloor}{\lceil kt/T \rceil - \lfloor kt/T \rfloor} (\eta_{\lceil
  kt/T \rceil} - \eta_{\lfloor kt/T \rfloor}),
\end{equation}
i.e., a linear interpolation between the keypoints.
\begin{itemize}
  \item[(a)] the metaparameter $\bm{\eta} \in \mathbb{R}^k$ is a
    vector of $k$ keypoints;
  \item[(b)] the algorithm $\mathcal{A}$ maps metaparameters
    $\bm{\eta}$ to a trained model $\mathcal{A}(\bm{\eta})$ by
    training a model for $T$ iterations with the learning rate schedule
    defined by \eqref{eq:lr_schedule};
  \item[(c)] the output function $\phi$ evaluates average loss on the
    validation set $\mathbf{X}_{val}$.
\end{itemize}

\subsubsection{Algorithm}
Following the theme of the rest of this section, we optimize the
metaparameter $\bm{\eta}$
directly using MGD. In particular, we initialize the keypoints to be
a flat learning rate schedule,
and then update the keypoints using the metagradient with respect to
the validation loss,
\begin{align*}
  \bm{\eta}^{(t+1)} = \bm{\eta}^{(t)} - \alpha \cdot
  \mathrm{sign}\left(\nabla_{\bm{\eta}}
  \phi(\mathcal{A}(\bm{\eta}^{(t)}))\right).
\end{align*}

\begin{figure}[tp]
    \centering
    \pgfplotsset{colormap/Set2}
\pgfplotsset{scaled y ticks=false}
\pgfplotsset{grid style={dashed,gray}}

\begin{tikzpicture}
    \begin{axis}[
      xlabel={Metagradient steps},
	  ylabel={Accuracy},
      height={5.2cm},
	  width={10cm},
      colormap name={Set2},
	  legend pos=south east,
      legend style={at={(1.1,1)}, anchor=north west, inner sep={6pt}},
      legend cell align={left},
	  xmin=-1,
	  xmax=51,
	  ymin=0.87,
	  ymax=0.95,
      grid=both
    ]
	\addlegendentry{Best grid target acc}
	\addplot[mark=none, index of colormap=3, domain=-2:1000, samples=2, dashed, very thick] {0.9401889};

	\addlegendentry{Best grid test acc}
	\addplot[mark=none, index of colormap=3, domain=-2:1000, samples=2, solid, very thick] {0.9266};

	\addplot[
      index of colormap=0,
	  mark size=0.0,
	  very thick,
	  dashed,
      error bars/.cd,
      y dir=both,
      y explicit,
    ] table[
      x=it,
      y=ymean,
      col sep=comma   %
	]{pgffigs/lr_opt_result/target_over_time.csv};
	\addlegendentry{MGD-LR target acc}

	\addplot[
      index of colormap=0,
	  mark size=0.0,
	  very thick,
      error bars/.cd,
      y dir=both,
      y explicit,
    ] table[
      x=it,
      y=ymean,
      col sep=comma   %
	]{pgffigs/lr_opt_result/test_over_time.csv};
	\addlegendentry{MGD-LR test acc}

	\addplot [
	  name path=upper_target,
	  draw=none,
	  forget plot
	] 
	table [
	  x=it, 
	  y expr=\thisrow{ymean} + \thisrow{yerr}, 
	  col sep=comma
	]{pgffigs/lr_opt_result/target_over_time.csv};

	\addplot [
	  name path=lower_target,
	  draw=none,
	  forget plot
	] 
	table [
	  x=it, 
	  y expr=\thisrow{ymean} - \thisrow{yerr}, 
	  col sep=comma
	]{pgffigs/lr_opt_result/target_over_time.csv};

    \addplot[
      fill opacity=0.2,
	  index of colormap=0,
    ] fill between[
      of=upper_target and lower_target,
    ];

	\addplot [
	  name path=upper_test,
	  draw=none,
	  forget plot
	] 
	table [
	  x=it, 
	  y expr=\thisrow{ymean} + \thisrow{yerr}, 
	  col sep=comma
	]{pgffigs/lr_opt_result/test_over_time.csv};

	\addplot [
	  name path=lower_test,
	  draw=none,
	  forget plot
	] 
	table [
	  x=it, 
	  y expr=\thisrow{ymean} - \thisrow{yerr}, 
	  col sep=comma
	]{pgffigs/lr_opt_result/test_over_time.csv};

    \addplot[
      fill opacity=0.2,
	  index of colormap=0,
    ] fill between[
      of=upper_test and lower_test,
    ];

    \end{axis}
\end{tikzpicture}

    \caption{Target and test accuracies of MGD's learning rate schedule
      over time closely match or exceed those found by a grid search over
      hundreds of combinations of hyperparameters.  95\% confidence
    intervals are plotted for MGD's results.}
    \label{fig:mgd_vs_grid}
\end{figure}

\subsubsection{Evaluation}
We aim to select the learning rate schedule that minimizes the
expected test set loss.
To do so, we reserve 90\% of the CIFAR-10 test set as a ``validation set''
on which we select hyperparameters.
We then use the remaining 10\% as a test set. We compare the
following two approaches:
\begin{itemize}
  \item {\bf Grid search:} We construct a grid over different one cycle
    learning rate schedules, varying the peak learning rate, starting learning
    rate, ending learning rate, and peak learning rate time. In total, we
    consider over 1,000 different learning rate schedules. We use the reserved
    90\% of the test set to select the best learning rate schedule
    from the grid.
  \item {\bf Metagradient descent (MGD):} We run 50 steps of
    MGD starting from a highly suboptimal flat learning rate schedule, aiming to
    minimize loss on the reserved 90\% of the test set. We use the
    last iteration
    of MGD as our learned learning rate schedule.
\end{itemize}
We evaluate the performance of each final learning rate schedule on the held-out
10\% test set and average the results over the same set of 5 unseen
random seeds.
 

\paragraph{Results.}
Comparing our learned hyperparameter schedule to grid search,
as shown in Figure \ref{fig:mgd_vs_grid}, our learned schedule
using only 50 steps of MGD matches the performance of the state-of-the-art
onecycle schedule found via grid search over more than 1000 configurations.
An important caveat, however, is that these numbers are not directly comparable:
grid search can be run in parallel across many machines, while steps
of MGD must be run sequentially.

In practice, we do not advise using MGD for optimizing low-dimensional
hyperparameters, especially ones that have been thoroughly optimized by grid
search (such as CIFAR-10 learning rate schedules
\citep{smith2017super,page2018cifar,li2019exponential,jordan202494}). Still, an
interesting avenue for future work is to study the utility of MGD for optimizing
high-dimensional hyperparameters that are less well-studied, such as
per-parameter/layer learning rates/weight decays for language models, attention
hyperparameters, or gradient preconditioners.
